% ch2_answers.tex -- Module 7, Chapter 2: answers to the Quick check questions.
\begin{chapteranswers}{2}{Light as a Wave}
  % ---- 2.1 Newton's and Huygens' models
  \item $v_\mathrm{x} = \dfrac{c}{n_\mathrm{x}} = \dfrac{\qty{3.00e8}{m.s^{-1}}}{2.42}
        = \qty{1.240e8}{m.s^{-1}}$. Answer: \final{$v = \qty{1.24e8}{m.s^{-1}}$}.
  \item Both models predict the same result. Newton: the corpuscle bounces like a
        perfectly bouncy ball. The part of its velocity along the surface stays the
        same. The part along the normal reverses. Huygens: reflected wavelets form a
        new wavefront that leaves at the same angle. Both give angle of incidence $=$
        angle of reflection. An experiment can only decide between two models if they
        predict \emph{different} results.
  \item Newton's model predicts that light is \emph{faster} in water than in air (the
        water pulls the corpuscles in). Huygens' model predicts that it is
        \emph{slower}. Foucault found that light is slower in water: the shift of the
        returning image was bigger with water in the path. This ruled out Newton's
        explanation of refraction and agreed with Huygens.
  % ---- 2.2 Diffraction
  \item Compare the gap with each wavelength: $\dfrac{2.0}{1.5} = 1.3$ and
        $\dfrac{2.0}{0.30} = 6.7$. For the \final{$\lambda = \qty{1.5}{cm}$} waves the gap
        is about one wavelength, so they spread out much more. For the \qty{0.30}{cm}
        waves the gap is almost 7 wavelengths wide, so most of the wave goes straight
        through.
  \item The spreading depends on $\lambda \div a$ ($a$ = slit width).
        (a) A smaller $a$ makes $\lambda \div a$ bigger, so the pattern (the central
        maximum and the bands beside it) becomes \final{wider}.
        (b) Blue light has a shorter $\lambda$ than red, so $\lambda \div a$ is smaller and
        the pattern becomes \final{narrower}.
  \item Particles that arrive at a place can only add more light. They cannot make a place
        darker. Waves from different parts of the slit can arrive \ang{180} out of phase:
        crest meets trough and they cancel. So the dark bands need waves, and they are
        evidence for the wave model.
  % ---- 2.3 Young's double slit
  \item The centre of the screen is the same distance from both slits, so the path
        difference is $\delta = 0 = 0\lambda$. This is true for every wavelength, so the
        waves always arrive in phase and interfere constructively: the centre ($m = 0$)
        is bright. With white light, every colour is bright at the centre at the same
        time, and all the colours together look white.
  \item Change to metres: $\lambda = 450 \times 10^{-9} = \qty{4.50e-7}{m}$ and
        $d = 0.200 \times 10^{-3} = \qty{2.00e-4}{m}$.
        $\Delta y = \dfrac{\lambda L}{d} = \dfrac{\qty{4.50e-7}{m} \times \qty{2.40}{m}}
        {\qty{2.00e-4}{m}} = \qty{5.40e-3}{m}$. Answer:
        \final{$\Delta y = \qty{5.40}{mm}$}. (The first fringe is at only \ang{0.13}, so
        the small-angle formula is fine.)
  \item Use $\Delta y = \dfrac{\lambda L}{d}$.
        (a) Green light has a shorter $\lambda$, so the fringes move \final{closer together}.
        (b) A bigger $L$ moves them \final{further apart}.
        (c) A bigger $d$ moves them \final{closer together}.
  \item $\lambda = \qty{6.00e-7}{m}$, $d = \qty{3.00e-4}{m}$, second dark fringe so $m = 2$
        and $m - \tfrac{1}{2} = 1.5$.
        $\sin\theta = \dfrac{\left(m - \frac{1}{2}\right)\lambda}{d}$ (divide both sides of
        $d\sin\theta = \left(m - \frac{1}{2}\right)\lambda$ by $d$)
        $= \dfrac{1.5 \times \qty{6.00e-7}{m}}{\qty{3.00e-4}{m}} = \num{3.00e-3}$, so
        \final{$\theta = \ang{0.172}$}.
        $y = L\tan\theta = \qty{2.00}{m} \times \tan\ang{0.172} = \qty{6.00e-3}{m}$:
        \final{$y = \qty{6.00}{mm}$}. Check: $\Delta y = \dfrac{\lambda L}{d} = \qty{4.00}{mm}$,
        and the second dark fringe is 1.5 spacings out: $1.5 \times 4.00 = \qty{6.00}{mm}$.
  \item (a) $v = f\lambda$, so $\lambda = \dfrac{v}{f}$ (divide both sides by $f$)
        $= \dfrac{\qty{340}{m.s^{-1}}}{\qty{850}{Hz}} = \qty{0.400}{m}$.
        Path difference: $\delta = \qty{3.60}{m} - \qty{3.00}{m} = \qty{0.60}{m}$, and
        $\dfrac{\delta}{\lambda} = \dfrac{\qty{0.60}{m}}{\qty{0.400}{m}} = 1.5$. One and a half
        wavelengths: crest meets trough, so the waves cancel (destructive). The sound is
        \final{quiet}.
        (b) Loud needs $\delta = m\lambda$ with $m$ a whole number. The lowest frequency has
        the longest wavelength, so take $m = 1$: $\lambda = \delta = \qty{0.60}{m}$. Then
        $f = \dfrac{v}{\lambda} = \dfrac{\qty{340}{m.s^{-1}}}{\qty{0.60}{m}} = \qty{567}{Hz}$.
        Answer: \final{$f \approx \qty{570}{Hz}$} (2 significant figures, like
        \qty{0.60}{m}). Check: at \qty{567}{Hz}, $\lambda = \qty{0.60}{m} = \delta$, a whole
        wavelength \cmark.
  % ---- 2.4 Diffraction gratings
  \item $N = 800 \times 1000 = \num{8.00e5}$ lines per metre, so
        $d = \dfrac{1}{N} = \dfrac{1}{\num{8.00e5}\,\unit{m^{-1}}}$:
        \final{$d = \qty{1.25e-6}{m}$}.
        First order: $\sin\theta = \dfrac{m\lambda}{d}$ (divide both sides of
        $d\sin\theta = m\lambda$ by $d$) $= \dfrac{1 \times \qty{4.45e-7}{m}}
        {\qty{1.25e-6}{m}} = 0.356$, so \final{$\theta = \ang{20.9}$}.
        Highest order: $\dfrac{d}{\lambda} = \dfrac{\num{1.25e-6}}{\num{4.45e-7}} = 2.81$.
        Round down: \final{$m = 2$} (at \ang{45.4}). $m = 3$ would need
        $\sin\theta = 1.07$, which is impossible.
  \item $N = 400 \times 1000 = \num{4.00e5}$ lines per m, so
        $d = \dfrac{1}{\num{4.00e5}\,\unit{m^{-1}}} = \qty{2.50e-6}{m}$.
        $y = \qty{0.640}{m} \div 2 = \qty{0.320}{m}$ (centre to one spot).
        $\tan\theta = \dfrac{y}{L} = \dfrac{\qty{0.320}{m}}{\qty{1.20}{m}} = 0.2667$, so
        $\theta = \tan^{-1}(0.2667) = \ang{14.93}$ and $\sin\theta = 0.2577$.
        $\lambda = \dfrac{d\sin\theta}{m} = \dfrac{\qty{2.50e-6}{m} \times 0.2577}{1}
        = \qty{6.44e-7}{m}$. Answer: \final{$\lambda = \qty{644}{nm}$} (red light).
  \item \final{Violet}. From $d\sin\theta = m\lambda$, $\sin\theta = \dfrac{m\lambda}{d}$.
        In the first order $m = 1$, and $d$ is the same for every colour. So the shortest
        wavelength (violet, about \qty{400}{nm}) has the smallest $\sin\theta$, and so
        the smallest angle from the centre. Red has the longest wavelength and goes
        furthest out.
  % ---- 2.5 Polarisation
  \item The light is already polarised, so use Malus' law:
        $I = I_\mathrm{max}\cos^2\theta = \qty{20.0}{W.m^{-2}} \times (\cos\ang{60.0})^2
        = 20.0 \times 0.250$. Answer: \final{$I = \qty{5.00}{W.m^{-2}}$}.
  \item A lamp's light is unpolarised, so one filter passes half of it at any angle. The
        \ang{30} is not needed. $I_0 = 2I = 2 \times \qty{12}{W.m^{-2}}$. Answer:
        \final{$I_0 = \qty{24}{W.m^{-2}}$}.
  \item $I = I_\mathrm{max}\cos^2\theta$, so $\cos^2\theta = \dfrac{I}{I_\mathrm{max}} = 0.40$
        (divide both sides by $I_\mathrm{max}$). Take the square root:
        $\cos\theta = \sqrt{0.40} = 0.6325$. Then $\theta = \cos^{-1}(0.6325)$:
        \final{$\theta = \ang{50.8}$}.
  \item First filter (unpolarised light): $I_1 = \tfrac{1}{2} \times \qty{200}{W.m^{-2}}
        = \qty{100}{W.m^{-2}}$. Second filter, \ang{20} from the first:
        $I_2 = 100 \times (\cos\ang{20})^2 = 100 \times 0.883 = \qty{88.3}{W.m^{-2}}$.
        Third filter, $\ang{90} - \ang{20} = \ang{70}$ from the second:
        $I_3 = 88.3 \times (\cos\ang{70})^2 = 88.3 \times 0.117 = \qty{10.3}{W.m^{-2}}$.
        Answer: \final{$I = \qty{10.3}{W.m^{-2}}$}.
  \item Sound is a longitudinal wave: the air vibrates back and forth along the direction
        of travel. There is no sideways direction of vibration for a filter to select, so
        sound cannot be polarised. Light is a transverse wave: its electric field
        vibrates at \ang{90} to the direction of travel, in many possible directions. A
        polariser can pass one of these directions and block the others.
\end{chapteranswers}
